% Part: counterfactuals
% Chapter: introduction
% Section: paradoxes-material

\documentclass[../../../include/open-logic-section]{subfiles}

\begin{document}

\olfileid[zh]{cnt}{int}{par}

\olsection{实质条件句的悖论}

C.~I. \zhFirst{Lewis}{刘易斯} 是最早批评用 $!A \lif !B$ 来符号化英语中「如果……那么……」陈述的人之一。\zhFirst{Lewis}{刘易斯} 所批评的是 Whitehead 和 \zhFirst{Russell}{罗素} 的\emph{Principia Mathematica}（《数学原理》）中实质条件句的用法，他们把 $\lif$ 读作「蕴含」。\zhFirst{Lewis}{刘易斯} 正确地抱怨道，如果 $\lif$ 表示「蕴含」，那么任何假的命题~$p$ 都蕴含 $p$ 蕴含~$q$（因为当~$p$ 为假时 $p \lif (p \lif q)$ 为真），并且任何真的命题~$q$ 都蕴含 $p$ 蕴含~$q$（因为当 $q$~为真时 $q \lif (p \lif q)$ 为真）。

逻辑学家当然知道，\emph{蕴含}，即逻辑衍推，并不是联结词，而是!!{formula}或语句之间的一种关系。所以为避免混淆，我们不应把 $\lif$ 读作「蕴含」。\footnote{尽管哲学家试图通过把它读作「仅当」来避免 \zhFirst{Lewis}{刘易斯} 所指出的混淆，但把「$\lif$」读作「蕴含」至今仍被数学家和计算机科学家广泛采用。} 只要我们不这样读，\zhFirst{Lewis}{刘易斯} 特有的那种担忧就不会出现：即使我们把 $p$ 当作一个碰巧为假的英语语句的代表，$p$ 也不「蕴含」$q$。要确定 $p \Entails q$，我们必须考虑\emph{全部} !!{valuation}；而且即使我们用 $p$ 来符号化一个碰巧为假的语句，$p \Entails/ q$ 仍然成立。

但「如果……那么……」这样的陈述仍有一些怪异之处，例如 \zhFirst{Lewis}{刘易斯} 的
\begin{quote}
如果月亮是由绿奶酪做成的，那么 $2+2=4$。
\end{quote}
以及如下推理：
\begin{quote}
  月亮不是由绿奶酪做成的。因此，如果月亮是由绿奶酪做成的，那么 $2+2=4$。

  $2+2 = 4$。因此，如果月亮是由绿奶酪做成的，那么 $2+2=4$。
\end{quote}
然而，如果「如果……那么……」就是 $\lif$，那么上面这个语句将毫无问题地为真，而这些推理也将毫无问题地有效。

另一个例子涉及重言式 $(!A \lif !B) \lor (!B \lif
!A)$。这暗示着，如果你从报纸上随机挑出两个直陈语句~$S$ 和 $T$，那么「如果 $S$ 则 $T$，或者如果 $T$ 则~$S$」这个语句应当为真。

\end{document}
